\section{Revisão da Literatura}
\label{sec:literatura}

Compreender a relação entre expansão agropecuária e conversão da vegetação nativa envolve discutir em quais condições os incentivos à produção favorecem a abertura de novas áreas ou o aproveitamento de terras já utilizadas. Essa questão aproxima o debate sobre retornos privados e custos ambientais das discussões sobre preços, capacidade produtiva e instituições. Para o problema proposto neste artigo, a revisão dessas contribuições busca esclarecer os mecanismos sugeridos pela literatura, as evidências que os sustentam e os limites de sua aplicação ao contexto de Mato Grosso. A trajetória econômica e ambiental do estado oferece o contexto para discutir a pertinência dessas explicações e suas possíveis diferenças entre atividades agropecuárias e biomas.

% Origem: introdução, REF1--REF2 e fundamentação conceitual de REF7/REF11.
\subsection{Expansão agropecuária, retornos privados e custos públicos}

A expansão agropecuária envolve escolha econômica entre usos alternativos da terra, na qual a conversão da vegetação em área produtiva para lavoura ou pecuária pode ampliar a renda privada associada à produção, ao mesmo tempo em que reduz estoques de capital natural e serviços ecossistêmicos que nem sempre são incorporados aos retornos percebidos pelos agentes. Essa diferença entre o valor social da vegetação e os benefícios que o proprietário consegue apropriar integra o enquadramento de \citeonline{Angelsen2010}, que distingue as rendas da agricultura e da floresta. A decisão depende, portanto, da comparação entre alternativas efetivamente acessíveis ao agente, e não apenas do valor da produção que poderá ser obtida na área convertida.

Nesse sentido, os efeitos da conversão podem ultrapassar os limites da propriedade convertida, embora a extensão desses efeitos dependa dos processos considerados. No Cerrado brasileiro, por exemplo, \citeonline{Schwamback2024} documentaram maior escoamento de nutrientes sob pastagem e maior perda de solo sob cana, em comparação com o cerrado arborizado. Na fronteira Amazônia-Cerrado, a substituição de floresta por soja esteve associada à redução do carbono do solo em até 100 cm de profundidade, com respostas diferentes conforme a profundidade e o manejo \cite{Popin2025}. Esses estudos documentam alterações biofísicas, mas não mensuram, por si sós, o custo social total da conversão.

Além das alterações biofísicas documentadas, a conversão pode ter implicações econômicas ao modificar serviços do solo que sustentam a produção ao longo do tempo. Nessa perspectiva, \citeonline{Brady2015} propuseram valorar esses serviços a partir de seus efeitos sobre fluxos futuros de renda agrícola, utilizando dados de experimentos de longo prazo na Europa e nos Estados Unidos que comparam a produtividade agrícola sob diferentes práticas de adubação e manejo do solo.

Com efeito, o problema permanece ainda menos resolvido quando se comparam benefícios e custos privados, como o lucro de um fazendeiro, e públicos, como os impactos gerados para outros agentes e territórios. Isso ocorre porque os serviços ambientais diferem entre lugares, operam em escalas distintas e muitos não possuem preços diretamente observáveis. Estudos que integram informações econômicas e biofísicas mostram a importância da localização e das alternativas de uso da terra para valorar esses serviços e orientar a conservação \cite{Crossman2009,Gascoigne2011}. Assim, a oposição entre conservar e produzir pode refletir diferenças entre retornos apropriados privadamente e custos distribuídos socialmente.

Para a Amazônia brasileira, \citeonline{Mann2012} estimaram esse descompasso ao comparar o retorno econômico potencial da conversão da floresta para a produção de soja com o valor monetário dos serviços ecossistêmicos mantidos pela vegetação nativa. Para tanto, combinaram estimativas do retorno da soja calculadas por um modelo de produtividade e pela diferença entre receitas e custos, com a valoração dos serviços ecossistêmicos por preços-sombra derivados de uma função de produção Cobb-Douglas, que incorporou a produtividade primária da vegetação como indicador desses serviços. Sob as hipóteses adotadas, a conclusão foi que o valor dos serviços supera a renda da soja em parcela expressiva da área analisada, e a omissão desses serviços reduz o bem-estar social relativamente ao ótimo modelado.

No caso do Pantanal, existem tentativas de mensurar economicamente os serviços ambientais, mas os próprios resultados evidenciam dificuldades de escala, cobertura e comparabilidade. \citeonline[p.~3]{SeidlMoraes2000}, ao recalcularem para a sub-região da Nhecolândia \footnote{Região do Pantanal sul-mato-grossense (Bacia do Alto Paraguai) caracterizada por florestas sazonais, pastagens nativas e lagoas alcalinas, com tradicional manejo da pecuária bovina extensiva.} estimativas globais com dados locais, obtiveram valor anual de aproximadamente US\$ 5,8 mil por hectare, em dólares de 1994, equivalente a cerca de metade daquele derivado dos parâmetros globais, diferença atribuída às condições biofísicas regionais. 

Mais recentemente, \citeonline[p.~7]{FairbrassEtAl2025}, ao estruturarem indicadores de capital natural para o Pantanal, buscaram reunir informações sobre a extensão, o estado de conservação e o valor econômico dos ecossistemas, os serviços ambientais oferecidos e as atividades humanas que deles se beneficiam. Dos 29 indicadores previstos para os ativos ecossistêmicos, encontraram dados que atendiam aos critérios do estudo para apenas 14, menos da metade do conjunto. Essa cobertura incompleta limitou a caracterização das relações entre a economia e o ambiente, dado que, embora houvesse informações sobre atividades com valor de mercado, como pecuária, soja, pesca e turismo, faltavam indicadores sobre a condição e o valor econômico dos ecossistemas que as sustentam, sobre parte das emissões geradas por essas atividades e acerca dos processos ecológicos essenciais, como a ciclagem de nutrientes e a manutenção da biodiversidade.

Sublinhe-se que as dificuldades de cobertura observadas no caso do Pantanal no parágrafo anterior também se estendem a outros biomas. \citeonline{CampoliStivali2026} reuniram estudos de valoração econômica da vegetação nativa nos seis biomas brasileiros (Amazônia, Cerrado, Pantanal, Mata Atlântica, Caatinga e Pampa) e padronizaram estimativas de custo social do desmatamento em dólares de 2025 por hectare. Ao final, os autores concluíram que os estudos diferem tanto em relação aos benefícios da vegetação incluídos na valoração econômica quanto nos métodos empregados. 

Como se percebe, as diferenças entre os valores estimados para os ganhos econômicos ambientais podem refletir tanto as características dos biomas como as diferenças no que cada estudo conseguiu mensurar. Em conjunto, os trabalhos discutidos mostram por que o retorno privado da conversão não representa, necessariamente, seu resultado econômico para a sociedade: a renda obtida com o novo uso da terra pode coexistir com perdas de serviços ambientais cujos custos não recaem integralmente sobre quem decide converter. Essa distinção fundamenta o recorte deste artigo, voltado a investigar como os incentivos privados se relacionam com a conversão da vegetação nativa. Cabe, portanto, examinar como os preços agropecuários e a capacidade produtiva da terra modificam a atratividade dos usos alternativos disponíveis aos produtores.

Nessa perspectiva, a contribuição deste artigo situa-se na análise das decisões que transformam o estoque ambiental, e não na estimação do valor monetário de seus serviços. Busca-se esclarecer como incentivos privados e condições institucionais se relacionam com a conversão para diferentes usos agropecuários. Essa delimitação aproxima a discussão sobre o descompasso entre retornos privados e custos sociais de uma questão empiricamente observável, sem supor que as lacunas de valoração serão resolvidas pelo exame de H1.

% Origem: REF3, REF6--REF8 e diálogo empírico de REF11.
\subsection{Preços agropecuários e incentivos à conversão}

No argumento deste artigo, os preços e a capacidade produtiva dos produtos agrícola ou pecuário representam dimensões observáveis dos incentivos ao uso agropecuário. Mantidos os custos e as demais condições, uma elevação de preço representa maior acréscimo de receita por hectare onde a produção potencial do uso beneficiado é maior, sendo isso justamente a relação econômica que motiva H1. Sua tradução em conversão efetiva, contudo, depende da vegetação disponível, dos custos de abertura da área, das alternativas em terras já convertidas e das restrições institucionais. Trata-se de uma previsão condicional a investigar, não de uma consequência necessária de qualquer aumento de produtividade.

Conforme suscitado logo no início desta seção, a conversão da vegetação nativa pode ser analisada como resultado de escolhas econômicas condicionadas por retornos relativos. \citeonline{Angelsen2010} organiza as políticas de redução do desmatamento em torno das rendas relativas da agricultura e da floresta, distinguindo intervenções que reduzem a atratividade da expansão, elevam os benefícios apropriáveis da conservação ou restringem diretamente a conversão. Esse enquadramento permite compreender por que a intensificação produtiva pode ampliar a pressão local de expansão em determinadas condições, sem presumir que produzir mais implique necessariamente converter mais terra.

A evidência de \citeonline{LubowskiEtAl2008}, para os Estados Unidos, relaciona os retornos líquidos dos usos alternativos e a qualidade da terra às decisões dos proprietários. Na Amazônia brasileira, \citeonline{Pfaff1999} identifica a relevância das características da terra e dos custos de transporte, advertindo que os efeitos estimados de políticas podem refletir a própria seleção de onde elas são implementadas. \citeonline{ChomitzThomas2003}, por sua vez, analisam conversão e intensidade pecuária em escala espacial detalhada, considerando precipitação, solos, infraestrutura, acesso a mercados e proteção territorial. Essas contribuições sustentam a consideração das condições produtivas e de acesso na análise de uso da terra. A heterogeneidade espacial que documentam não equivale, por si só, a um teste da resposta diferencial a uma mesma variação de preços.

A distinção entre capacidade produtiva e produtividade observada é pertinente a essa passagem. O rendimento obtido em áreas já cultivadas também reflete decisões de manejo e utilização de insumos, além das características da terra selecionada para produzir. Uma medida de capacidade produtiva predeterminada precisa, portanto, ser avaliada quanto ao que representa e à sua posição temporal em relação à conversão. Da mesma forma, preços e produtividade, isoladamente, não revelam o benefício privado líquido esperado, pois não incorporam todos os custos, restrições e expectativas que integram a decisão do produtor.

Em escala tropical, \citeonline{BermanEtAl2023} combinaram preços internacionais e aptidão agrícola para examinar a relação entre incentivos de mercado e desmatamento entre 2001 e 2018.\footnote{A pesquisa abrangeu regiões tropicais de diferentes países, com destaque para a Amazônia, o Sudeste Asiático e a Bacia do Congo, onde a associação entre preços agrícolas e desmatamento foi mais forte.} Os autores encontraram uma associação positiva entre preços agrícolas e desmatamento, cuja intensidade variou conforme as características econômicas locais. Essa abordagem é pertinente à H1 porque permite considerar que um mesmo aumento de preço pode gerar incentivos distintos à conversão, dependendo das possibilidades produtivas de cada território. O uso da aptidão agrícola para medir a exposição aos preços, contudo, não demonstra, por si só, que a capacidade produtiva local modifica a resposta da conversão da forma prevista nesta pesquisa. Há também um desafio de identificação: os preços internacionais podem refletir choques de oferta, e é necessário justificar por que a exposição a suas variações pode ser considerada exógena ao desmatamento analisado. Portanto, o uso de preços internacionais não basta para assegurar uma interpretação causal dos resultados.

Sem embargo, a relação entre preços e conversão depende do uso que se expande e daquele que é substituído. Na bacia do Tapajós, \citeonline{Braganca2018} encontrou que a elevação do preço relativo de lavouras em relação à carne favoreceu a substituição de pastagens por cultivos (\textit{pasture-to-cropland}) e reduziu o desmatamento local. Uma mudança nos preços relativos entre atividades, portanto, não se confunde com uma elevação do retorno esperado da conversão da vegetação nativa. A pressão para abrir novas áreas agropastoris deve, assim, ser distinguida da reorganização da produção em áreas já convertidas. Para este artigo, essas margens de ajuste devem também ser distinguidas da intensificação produtiva. Soja, milho e pecuária tampouco podem ser combinados em um único índice sem explicitar unidades, exposição e mecanismo econômico.

À luz desses antecedentes, a contribuição aqui pretendida consiste em examinar as condições territoriais em que os incentivos de preços se associam à conversão da vegetação nativa, avançando sobre a relação documentada por \citeonline{BermanEtAl2023}, \citeonline{AssuncaoEtAl2015} e \citeonline{HardingEtAl2021}, dentre outros. O diferencial do artigo está em analisar conjuntamente a capacidade produtiva local, o ambiente institucional e os destinos da conversão, distinguindo a abertura de áreas para agricultura e pastagem da reorganização de usos em áreas anteriormente convertidas, destacada por \citeonline{Braganca2018}. Essa abordagem busca esclarecer por que movimentos semelhantes de preços podem se traduzir em respostas distintas de conversão e, com isso, oferecer evidências sobre o alcance e os limites das explicações existentes. Para o caso do estado de Mato Grosso, a contribuição pretendida é identificar quais atividades e territórios apresentam maior sensibilidade aos preços e em que medida essa resposta varia com as condições produtivas e institucionais, oferecendo subsídios à discussão de políticas de conservação que considerem essas diferenças e abrindo perspectivas para que estudos futuros investiguem a adequação territorial de instrumentos regulatórios e econômicos de conservação em outros contextos, regiões e biomas.

% Origem: REF1--REF2 e REF9; limites de escopo estabelecidos ao final de REF11.
\subsection{Instituições e respostas aos incentivos}

As instituições ajudam a compreender por que os produtores podem responder de formas distintas aos incentivos de preços discutidos anteriormente. Na economia institucional, \citeonline{North1991} compreende as instituições como regras formais, a exemplo de leis, regulamentos e direitos de propriedade, e restrições informais, como costumes e normas de conduta compartilhados por membros de uma comunidade. \citeonline{North1993} distingue essas regras das organizações que atuam sob elas e participam de sua elaboração e aplicação, entre as quais órgãos públicos, empresas e associações. Na governança dos recursos naturais, \citeonline{Ostrom2009} destaca a relevância das regras efetivamente praticadas, do monitoramento e das sanções, inclusive quando organizados pelos próprios usuários. Por seu turno, \citeonline{Shimshack2014} examina como o monitoramento e a aplicação das regras influenciam o cumprimento das exigências ambientais. Essas contribuições permitem situar as respostas aos incentivos na articulação entre normas, organizações e condições de implementação.

Neste artigo, o ambiente institucional compreende as regras formais e informais que orientam as decisões de conversão, conservação e uso dos recursos ambientais, bem como os mecanismos que sustentam sua aplicação. Inclui leis e regulamentos sobre o uso da terra e da vegetação, direitos de propriedade, costumes e normas de conduta reconhecidos pelos agentes. Sua implementação envolve organizações públicas e privadas e os próprios usuários dos recursos, por meio de monitoramento, fiscalização e aplicação de sanções legais, contratuais ou sociais. Para H1, importa compreender como essa articulação modifica os custos, os riscos e os retornos esperados das alternativas de uso da terra, distinguindo a existência de uma regra das condições que tornam seu cumprimento efetivo.

Sob essa perspectiva, o enquadramento econômico não supõe efeitos uniformes das instituições. \citeonline{PachecoMeyer2022} comparam regimes fundiários em diferentes contextos brasileiros e mostram que os resultados dependem do regime de origem, da alternativa de destinação da terra e do contexto socioambiental. A comparação de terras privadas com terras públicas sem destinação não produz a mesma conclusão que sua comparação com áreas protegidas. Já \citeonline{ProbstEtAl2020}, ao examinar o programa Terra Legal, encontram aumento do desmatamento entre pequenos e médios proprietários titulados e maior resposta aos preços agropecuários. A regularização, desse modo, não deve ser confundida com um indicador inequívoco de conservação.

Os estudos de \citeonline{PachecoMeyer2022} e \citeonline{ProbstEtAl2020} ajudam a compreender como os direitos de propriedade podem modificar os incentivos à conversão da vegetação, embora examinem questões distintas: o primeiro compara regimes fundiários, enquanto o segundo avalia uma iniciativa de titulação da terra. Para o presente artigo, essas evidências contextualizam a relação entre regras de acesso e uso da terra e a atratividade econômica da conversão, sem pressupor que a formalização de direitos favoreça necessariamente a conservação. A análise dos mecanismos de apropriação irregular e valorização especulativa da terra, como ponderado na introduçãoo, permanece reservada a outro artigo.

A mesma cautela se aplica ao Cadastro Ambiental Rural. \citeonline{JungEtAl2022} analisam Pará e Mato Grosso e identificam respostas distintas da conversão para lavouras e pastagens, com heterogeneidade por porte dos estabelecimentos e conformidade ambiental. Em seu modelo dinâmico de uso da terra, com procedimentos para enfrentar a seleção no cadastro, não encontram efeito estatisticamente significativo da inscrição sobre a conversão para lavouras, enquanto identificam respostas da expansão de pastagens que variam entre os grupos examinados. Os resultados apontam o acesso ao crédito como um possível mecanismo. 

Essa evidência reforça a necessidade de examinar como instrumentos distintos afetam oportunidades produtivas e custos de conversão, em vez de reunir todo registro administrativo sob uma medida única de qualidade institucional. O resultado não significativo para lavouras não estabelece que o efeito seja exatamente zero nem permite concluir que o cadastro seja irrelevante para qualquer uso da terra. Em conjunto com os demais resultados do estudo, ele mostra por que o destino da conversão e as condições dos produtores importam para interpretar o alcance de uma intervenção.

Fora do contexto brasileiro, \citeonline{ViallonEtAl2019} mostram, em doze casos de uso da terra na Suíça, que a implementação local depende da capacidade das autoridades e dos direitos adquiridos pelos proprietários. \citeonline{AnderssonEtAl2014}, a partir de comunidades florestais na Bolívia, relacionam condições das florestas a diferentes combinações de elaboração de regras, monitoramento e sanções organizadas pelos usuários. Esses trabalhos dão conteúdo empírico à distinção entre as regras e sua aplicação, incluindo a atuação dos próprios agentes na organização do uso dos recursos. A pertinência desses mecanismos em contextos específicos motiva a atenção às condições locais, sem estabelecer a superioridade de um arranjo de governança para Mato Grosso.

No âmbito privado, a Moratória da Soja, examinada por \citeonline{HardingEtAl2021}, exemplifica a implementação de critérios de conservação por meio das relações de compra na cadeia produtiva. Esse caso situa os critérios comerciais entre os mecanismos que podem alterar as alternativas dos produtores. Instrumentos econômicos de conservação, como pagamentos por serviços ambientais, mercados de carbono ou financiamento ambiental, também integram esse contexto. As receitas potenciais associadas à conservação do carbono florestal, por exemplo, podem modificar os retornos relativos enfrentados pelos agentes \cite{Angelsen2010,SticklerEtAl2013}. Esses exemplos contextualizam a diversidade de incentivos à conservação; a avaliação direta da eficácia desses instrumentos permanece fora do recorte empírico deste artigo.

Em Mato Grosso, esse enquadramento orienta o exame da sensibilidade da conversão aos preços em contextos institucionais distintos, considerando a capacidade produtiva local e os destinos agropecuários da área convertida. Conforme delimitado na introdução, os registros de autuações do IBAMA constituem uma fonte candidata para observar uma dimensão da atuação fiscalizatória. Esses registros não representam, isoladamente, a efetividade das regras nem permitem mensurar todo o ambiente institucional, incluindo costumes e normas informais. A contribuição pretendida está em relacionar as dimensões institucionais que os dados permitirem observar às diferenças na resposta da conversão, mantendo a definição conceitual mais ampla como contexto de interpretação. Cabe, a seguir, situar essas relações na trajetória econômica e ambiental de Mato Grosso desde os anos 2000.

% Origem: REF2, REF9 e comparação entre instrumentos em REF11.
% Origem: REF4--REF6, REF10 e contribuição/transição empírica de REF11.
\subsection{Situação econômica e ambiental de Mato Grosso}

A situação econômica e ambiental de Mato Grosso reúne os processos discutidos nas subseções anteriores: a expansão agropecuária transformou o estoque de vegetação nativa por ter sido o fator terra determinante para a expansão agropecuária do estado, enquanto as instituições condicionaram os usos da terra. 

A integração aos mercados internacionais constitui parte desse contexto. O ingresso da China na Organização Mundial do Comércio, em dezembro de 2001, marcou uma etapa de sua inserção no comércio mundial \cite{WTO2021}. Em Mato Grosso, \citeonline{LathuilliereEtAl2014} documentam que, entre 2006 e 2010, a China superou a Europa como destino das exportações de soja, acompanhando mudanças na escala e nas implicações ambientais da produção. Essa reorientação ajuda a situar a importância da demanda externa para os usos agropecuários da terra. A transmissão desses incentivos às decisões locais, contudo, depende dos preços recebidos, do câmbio, dos custos e das condições de acesso aos mercados, sem que a coincidência temporal permita atribuir a expansão estadual exclusivamente à adesão chinesa à OMC.

A dimensão desse crescimento aparece nas estatísticas oficiais. Entre 2000 e 2024, a produção de soja em Mato Grosso passou de aproximadamente 8,8 para 38,4 milhões de toneladas, enquanto a de milho aumentou de 1,4 para 47,4 milhões de toneladas \cite{IBGEPAM2026}. No mesmo intervalo, o rebanho bovino passou de 18,9 para 32,9 milhões de cabeças \cite{IBGEPPM2026}. Esses números situam a escala das atividades examinadas, mas o aumento da produção ou do efetivo animal não informa, isoladamente, quanto de vegetação foi convertido. É necessário distinguir a ampliação da área utilizada, a substituição entre usos e o aumento da produção nas áreas já ocupadas.

Parte dessa transformação foi viabilizada pela difusão de tecnologias adaptadas à agricultura tropical. \citeonline{Spera2017} destaca a adaptação de cultivares de soja, a correção dos solos e o desenvolvimento de variedades de ciclo curto, que viabilizam o cultivo sequencial de soja e milho no mesmo ano agrícola. Em Mato Grosso, \citeonline{VanWeyEtAl2013} mostram que a expansão de duas safras está associada ao acesso a transportes, à produção agrícola anterior e à conexão com mercados. Esses processos ampliam as alternativas de produção em terras já convertidas e podem modificar o retorno da expansão agropecuária. Para H1, reforçam a necessidade de distinguir a capacidade produtiva local predeterminada do rendimento observado, que também incorpora escolhas de manejo, insumos e tecnologia.

A trajetória ambiental não acompanhou a produção de maneira constante. Na porção florestal de Mato Grosso, \citeonline{MacedoEtAl2012} verificaram que 91\% da expansão da soja entre 2006 e 2010 ocorreu em terras anteriormente abertas, sobretudo pastagens, ao mesmo tempo em que a produção cresceu e o desmatamento diminuiu. A conversão para pastagem respondeu pela maior parte do desmatamento remanescente examinado. Esse resultado mostra que crescimento produtivo e menor abertura de florestas podem coexistir em determinadas condições, sem estabelecer que a intensificação assegure a conservação. A mudança entre períodos, em diálogo com a substituição de usos discutida por \citeonline{Braganca2018}, torna relevante examinar quando os incentivos se traduzem em novas conversões e quando favorecem ajustes em áreas já utilizadas.

Os efeitos econômicos dessa expansão ultrapassaram a atividade rural. \citeonline{RichardsEtAl2015} relacionam a produção e a rentabilidade da soja ao crescimento do emprego e da renda gerada em atividades não agrícolas em Mato Grosso, por meio da demanda por serviços, bens e investimentos. O estudo também destaca a dependência desses resultados em relação aos preços e ao câmbio. Há, portanto, razões para considerar que a importância econômica da agropecuária envolve tanto a renda apropriada pelos produtores quanto suas relações com a economia regional, preservando a distinção, estabelecida anteriormente, entre esses benefícios e os custos ambientais da conversão.

A projeção política de segmentos do setor constitui outra dimensão dessa transformação. Em análise qualitativa centrada em Mato Grosso, \citeonline{Ioris2016} discute os sinais de enriquecimento de produtores de soja e sua articulação com organizações representativas, como a Aprosoja-MT, e o poder público. Para o argumento institucional deste artigo, importa reconhecer que os agentes respondem às regras e também participam das disputas sobre seu conteúdo e sua implementação.

No campo ambiental, o Plano de Ação para Prevenção e Controle do Desmatamento na Amazônia Legal (PPCDAm), lançado em 2004, e o uso de alertas de satélite pelo sistema DETER fortaleceram a articulação entre monitoramento e fiscalização \cite{AssuncaoEtAl2023}. A Moratória da Soja, iniciada em 2006, acrescentou uma restrição comercial à expansão: empresas signatárias se comprometeram a não comprar soja produzida em áreas do bioma Amazônia desmatadas após a data de corte do acordo, posteriormente ajustada para 22 de julho de 2008. \citeonline{HeilmayrEtAl2020} estimam redução do desmatamento associada à Moratória em áreas aptas à soja e destacam sua complementaridade com cadastros públicos e monitoramento. Sua abrangência territorial também importa, pois os mesmos critérios não se estendiam ao Cerrado e ao Pantanal, tornando inadequado supor uma exposição institucional uniforme no estado.

O novo Código Florestal, instituído pela Lei n.~12.651/2012, reorganizou as regras de proteção da vegetação e estabeleceu o Cadastro Ambiental Rural (CAR) como registro nacional obrigatório dos imóveis rurais. A lei prevê percentuais distintos de Reserva Legal em alguns biomas brasieiros conforme a formação vegetal e instrumentos de regularização de passivos, de modo que as exigências não podem ser representadas por uma única proporção para todo o estado \cite{Brasil2012}.\footnote{O art.~12 estabelece, como regra geral na Amazônia Legal, 80\% para imóveis em áreas de florestas, 35\% em áreas de cerrado e 20\% em campos gerais, observadas as exceções legais. A Amazônia Legal é um recorte territorial distinto do bioma Amazônia.} \citeonline{SticklerEtAl2013} analisam os custos e as consequências potenciais do cumprimento dessas obrigações em Mato Grosso, discutindo a instabilidade regulatória e sua combinação com incentivos positivos à conservação. Como indica a evidência de \citeonline{JungEtAl2022}, discutida anteriormente, a inscrição no CAR tampouco permite presumir uma resposta homogênea de conservação entre atividades e produtores.

Os compromissos internacionais introduziram outras possibilidades de articulação entre produção e conservação. O Acordo de Paris, de 2015, reconheceu o papel das florestas na mitigação climática e incentiva o apoio à redução das emissões por desmatamento, inclusive por pagamentos por resultados \cite{UNFCCC2015}. No mesmo ano, Mato Grosso apresentou na COP21 a estratégia Produzir, Conservar e Incluir (PCI); em 2017, formalizou acordos com Alemanha e Reino Unido no âmbito do programa REDD Early Movers (REM), voltados ao financiamento de ações de conservação e desenvolvimento sustentável \cite{FagundesDomingues2017}. Esses marcos situam incentivos potenciais à manutenção da vegetação, cuja existência não equivale à demonstração de seus efeitos sobre as decisões dos produtores.

A evidência sobre implementação ajuda a interpretar essa diversidade de instrumentos. \citeonline{HardingEtAl2021} encontram, na Amazônia brasileira, respostas distintas do desmatamento aos preços conforme as políticas e as áreas consideradas, incluindo resultados compatíveis com substituições de cultivos e deslocamentos da pressão para terras não protegidas. \citeonline{AssuncaoEtAl2023}, ao explorar limitações de detecção por nuvens no DETER, identificam efeito de contenção do desmatamento pelo monitoramento e pela fiscalização. Em Mato Grosso, \citeonline{Sousa2016} examina multas, embargos e apreensões e destaca diferenças entre sanções e dificuldades de implementação. Esses resultados motivam o exame do condicionamento institucional da resposta aos preços, mas a contagem de autuações exige cautela: a fiscalização pode se concentrar justamente onde a pressão de conversão é maior e não se confunde com a efetividade das regras.

As diferenças territoriais completam esse quadro. Os mapas de \citeonline{SimoesEtAl2020} distinguem trajetórias de agricultura, pastagem e vegetação em Mato Grosso entre 2001 e 2017, abrangendo Amazônia, Cerrado e Pantanal. No Pantanal brasileiro, \citeonline{SeidlEtAl2001} examinam indicadores econômicos e incentivos dos proprietários relacionados à expansão de pastagens cultivadas, evidenciando uma dinâmica que não se resume à expansão de grãos. Sua abrangência inclui áreas fora de Mato Grosso, o que limita a transposição direta dos resultados. A presença dos três biomas oferece, por isso, uma oportunidade de comparar combinações de capacidade produtiva, usos da terra e instituições, desde que as diferenças biofísicas e de cobertura dos dados sejam consideradas.

À luz desse contexto, o interesse de Mato Grosso para a pesquisa decorre da combinação entre escala agropecuária, mudanças nas alternativas produtivas e condições institucionais distintas. A contribuição pretendida consiste em examinar, sob uma base territorial comum, se a resposta da conversão bruta aos preços varia com a capacidade produtiva local e com as dimensões institucionais observáveis, distinguindo os destinos agricultura e pastagem. Essa comparação procura esclarecer o alcance das relações previstas na hipótese princiap deste trabalho entre atividades e biomas, considerando também resultados nulos ou contrários às previsões e a precisão com que forem estimados. A conversão presente modifica a vegetação disponível para decisões posteriores, dando conteúdo ao tratamento do estoque ambiental como componente endógeno do processo econômico.
